\chapter{Conclusiones}
El desarrollo permitió abordar el diseño, cálculo analítico, modelado por simulación y validación experimental de
sistemas analógicos de acondicionamiento de señal y filtrado activo de orden superior. A partir de los resultados
obtenidos y de las dificultades operativas resueltas en el laboratorio, se formulan las siguientes conclusiones
técnicas:

En primer término, respecto al módulo de acondicionamiento para lazo de corriente de $4 \unit{\milli\ampere}$ a
$20 \unit{\milli\ampere}$, se corroboró que la estructura en tres etapas entrega un comportamiento estático lineal
con errores relativos sistemáticos inferiores al $-0.6\%$. Dicho desvío se atribuyó con exactitud a la corriente de
base del transistor bipolar NPN BC337 ($\alpha \approx \beta/(\beta+1)$) y a las tolerancias del $1\%$ en los
resistores de shunt.

En segundo término, el filtro pasa-bajos Sallen-Key de $10 \unit{\hertz}$ demostró un gran rechazo frente al ruido
electromagnético de red eléctrica: la señal de interferencia de $50 \unit{\hertz}$ se atenuó en $24.50 \unit{\decibel}$,
asegurando la inmunidad del lazo de corriente frente a acoplamientos parásitos de cables de potencia.

En tercer término, el filtro pasa-bajos de cuarto orden con aproximación Butterworth y topología Sallen-Key cumplió con
la distribución simétrica de ganancia en dos etapas ($K_1 = 3, K_2 = 3$), alcanzando los $19.08 \unit{\decibel}$ sin
requerir etapas amplificadoras adicionales. Las mediciones en osciloscopio dieron una atenuación de $0.56
\unit{\decibel}$ en el borde de banda de paso y una pendiente de transición abrupta superior a $-75
\unit{\decibel\per\decade}$.

En cuarto término, el filtro pasa-altos de cuarto orden en topología MFB con aproximación Chebyshev verificó la
efectividad de la técnica de polarización unipolar con masa virtual a $+12 \unit{\volt}$. El circuito entregó una
ganancia estable de $6.13 \unit{\volt\per\volt}$ en la banda de paso profunda y una atenuación relativa de $-3.30
\unit{\decibel}$ en torno al borde de transición de $10.87 \unit{\kilo\hertz}$, confirmando la marcada selectividad de
4to orden y demostrando la robustez de la estructura de realimentación múltiple frente a dispersiones capacitivas
comerciales.

Finalmente, el diseño analítico y la simulación del filtro pasa-banda de segundo orden en topología Sallen-Key
permitió validar la técnica de descomposición en cascada de etapas pasa-altos y pasa-bajos, centrando la resonancia en
$25.12 \unit{\kilo\hertz}$ con una selectividad coherente con el ancho de banda asignado de $2.0 \unit{\kilo\hertz}$.

En definitiva, la experiencia evidenció que el éxito en la implementación de filtros activos analógicos no reside
únicamente en el rigor algebraico de la síntesis de polinomios, sino de manera preponderante en la consideración de
los efectos no ideales de los componentes reales: producto ganancia-ancho de banda, impedancia de fuentes de corriente
internas, capacidades de sumidero a potencial de masa y tolerancias de las series comerciales normalizadas.
